\seccionreporte{Modelo de aprendizaje no supervisado}{sec:modelo}

\propositoseccion{%
  Documenta el algoritmo, su fundamentación breve y la evidencia de entrenamiento (notebook + métricas).%
}

\subsection{Selección del algoritmo}

Algoritmo implementado: \textbf{HDBSCAN + UMAP}.

Nos fuimos por HDBSCAN porque los proyectos no forman grupitos perfectos redondos como los que busca K-Means. Aquí tenemos densidades raras y proyectos que de plano no encajan con nada (outliers o Océanos Azules). HDBSCAN es buenísimo encontrando clústeres basados en la densidad de los datos sin que le tengamos que decir cuántos clústeres queremos desde antes.

\subsection{Fundamento breve}

Básicamente, lo que hace el pipeline completo es: primero, agarramos los textos y usamos un modelo de embeddings de lenguaje (SentenceTransformers) para convertirlos en vectores matemáticos de 384 dimensiones. Luego, metemos esos vectores gigantescos a UMAP para aplastarlos a solo 15 dimensiones, porque si le damos las 384 dimensiones directas a HDBSCAN, se pierde por la maldición de la dimensionalidad.

Una vez que UMAP reduce el espacio cuidando la topología, HDBSCAN busca dónde se agrupan los puntos. Los puntos que quedan volando solos, HDBSCAN los marca como ruido (con un -1), lo cual para nosotros significa que ese proyecto es súper original y no se parece a nada anterior.

\subsection{Entrenamiento y artefacto}

\begin{itemize}
  \item Notebook de evidencia: \texttt{notebooks/entrenamiento\_hdbscan.ipynb} (o el script equivalente \texttt{run\_training.py} en el repositorio).
  \item Semilla aleatoria: \texttt{42}, para que siempre nos dé los mismos resultados si corremos el script.
  \item Corpus: Usamos un conjunto sintético balanceado de 50 propuestas distribuidas en 10 dominios.
  \item Artefactos generados: \texttt{app/models/umap\_50d\_model.joblib} (88.4 KB) y \texttt{app/models/hdbscan\_model.joblib} (27.1 KB).
\end{itemize}

\subsection{Métricas de evaluación (no supervisado)}

Al correr el notebook, las métricas del modelo salieron bastante sólidas:

\begin{itemize}
  \item \textbf{Silhouette Score:} \textbf{0.6962}. Un score mayor a 0.5 ya es bueno, así que casi 0.7 significa que los clústeres quedaron súper bien separados unos de otros.
  \item \textbf{Davies-Bouldin:} \textbf{0.3411}. Como esta métrica penaliza la dispersión, entre más baja mejor. Un 0.34 indica clústeres muy compactitos.
  \item \textbf{Clústeres detectados:} \textbf{6} agrupaciones naturales (más los candidatos a ruido).
\end{itemize}

\figuraopcional{figures/mapa_semantico_2d.png}{1\textwidth}{%
  El mapa semántico generado en 2D por UMAP muestra claramente cómo se agrupan los dominios del corpus.%
}{Métricas o gráficos del entrenamiento}
